%% ma: L12-B08-0034
%% lop: 12
%% chuong: 2
%% bai: 8
%% dang: 12.8.27
%% loai: TN4
%% muc: NB
%% dap_an: "A"
%% nguon: "LaTeX giáo viên gửi 21/10/2026 (M05, Vectơ và hệ trục tọa độ trong không gian, 107 câu), Phần I Câu 10"
%% kien_thuc: "Bài 8 (tọa độ trọng tâm tứ diện)"
%% trang_thai: cho-duyet
%% ly_do: "Dạng mới đề xuất 12.8.27"
\begin{ex}
Trong không gian $Oxyz$, cho tứ diện $ABCD$ với $A(1;2;3)$, $B(0;1;1)$, $C(-1;2;0)$ và $D(4;3;0)$. Tọa độ trọng tâm $G$ của tứ diện là
\choice{\True $(1;2;1)$}{$(4;8;4)$}{$(2;4;2)$}{$(1;1;2)$}
\loigiai{$G\left(\dfrac{1+0-1+4}4;\dfrac{2+1+2+3}4;\dfrac{3+1+0+0}4\right)=(1;2;1)$.}
\end{ex}
